% Author: Isaac H. Lopez Diaz <isaac.lopez@upr.edu>
% Description: Phase 3 of Project CCOM 6050

\documentclass{article}

\usepackage{tabularx}

\title{Phase 3 \\ CCOM 6050: Design and Analysis of Algorithms}
\author{Isaac H. Lopez Diaz}

\begin{document}
\maketitle

% SECTION 1 %
\section{Statement of the Problem}
This report presents how a reflective language - a language able to reason about itself - can
benefit from a type system, specifically dependent types - types that depend on their value (term).
Most, if not all, reflective languages are extensions of some Lisp dialect, meaning that they are
likely to be dynamically-typed, so one can argue that there is a need for a reflective language
that incorporates a type system. One notable argument for incorporating type systems on a
reflective language is that one can ensure semantics-preserving rewrites (safe transformations).

% SECTION 2 %
\section{Background and Methodologies}
\subsection*{Background}
\subsection*{Methodologies}

% SECTION 3 %
\section{Algorithms}
\subsection*{Coquand's type-checking algorithm}
\subsection*{Partial Evaluation}

% SECTION 4 %
\section{Complexity Analysis}
\subsection*{Type-checking Algorithm}
\begin{tabularx}{0.8\textwidth}{
  | >{\raggedright\arraybackslash}X
  | >{\centering\arraybackslash}X
  | >{\raggedright\arraybackslash}X
  | >{\raggedright\arraybackslash}X | }
  \hline
  Case & Time Complexity & Space Complexity \\
  \hline
  Best Case & $\Theta(1)$ & $\Theta(1)$ \\
  \hline
  Average Case & $\Theta(n)$ & $\Theta(n)$ \\
  \hline
  Worst Case & $\Theta(n^2)$ or $\Theta(n^3)$ & $\Theta(n^2)$ \\
  \hline
\end{tabularx}

\subsection*{Partial Evaluation Algorithm}
\begin{tabularx}{0.8\textwidth}{
  | >{\raggedright\arraybackslash}X
  | >{\centering\arraybackslash}X
  | >{\raggedright\arraybackslash}X
  | >{\raggedright\arraybackslash}X | }
  \hline
  Case & Time Complexity & Space Complexity \\
  \hline
  Best Case & $\Theta(1)$ & $\Theta(1)$ \\
  \hline
  Average Case & $\Theta(n)$ & $\Theta(n)$ \\
  \hline
  Worst Case & $\Theta(n^2)$ or $\Theta(n^3)$ & $\Theta(n^2)$ or $\Theta(n^3)$ \\
  \hline
\end{tabularx}

% SECTION 5 %
\section{Computational Tools}

% BIBLIOGRAPHY %
\bibliography{refs}
\nocite*
\bibliographystyle{acm}
\end{document}
